\section{AMC 12 12B 2014}
        \begin{enumerate}
        	\item (\textit{AMC 12 12B 2014}) Leah has $13$ coins, all of which are pennies and nickels. If she had one more nickel than she has now, then she would have the same number of pennies and nickels. In cents, how much are Leah's coins worth?


 $\textbf{(A)}\ 33\qquad\textbf{(B)}\ 35\qquad\textbf{(C)}\ 37\qquad\textbf{(D)}\ 39\qquad\textbf{(E)}\ 41$

 \textbf{Answer: }C



\textbf{Solution}

She has $p$ pennies and $n$ nickels, where $n + p = 13$. If she had $n+1$ nickels then $n+1 = p$, so $2n+ 1 = 13$ and $n=6$. So she has 6 nickels and 7 pennies, which clearly have a value of $\boxed{\textbf{(C)}\ 37}$ cents.

 


	\item (\textit{AMC 12 12B 2014}) Orvin went to the store with just enough money to buy $30$ balloons. When he arrived he discovered that the store had a special sale on balloons: buy $1$ balloon at the regular price and get a second at $\frac{1}{3}$ off the regular price. What is the greatest number of balloons Orvin could buy?


 $\textbf{(A)}\ 33\qquad\textbf{(B)}\ 34\qquad\textbf{(C)}\ 36\qquad\textbf{(D)}\ 38\qquad\textbf{(E)}\ 39$

 \textbf{Answer: }C



\textbf{Solution 1 (Rigorous)}

Let us assume that Orvin will buy $a$ balloons at the original price $k$, and $b$ balloons at the discounted price.

 
$ak+\frac{2b}{3}k\leq{30k}$
$a+\frac{2b}{3}\leq{30}$

 
Because the number of balloons with discount cannot exceed the number of balloons with the original price, $b\leq{a}$.

 
$b+\frac{2b}{3}\leq a+\frac{2b}{3}\leq{30}$
$\frac{5b}{3}\leq{30}$
$b\leq{18}$

 
Let $Q=a+b$ be the total number of balloons Orvin buys.

 
$a+\frac{2b}{3}=Q-\frac{b}{3}\leq{30}$
$Q\leq{30+\frac{b}{3}}$

 
Substitute in $b\leq{18}$.

 $Q\leq{30+\frac{b}{3}}\leq{30+\frac{18}{3}}=36$

 
The maximum value of $Q$ is indeed attainable. Buying $18$ balloons at the original price and $18$ balloons at the discounted price, Orvin will spend $18k+18\times{\frac{2}{3}}k=30k$ dollars. This completes the proof.

 
The answer to the problem is $36$.

 


 



\textbf{Solution 2}

If every balloon costs $n$ dollars, then Orvin has $30n$ dollars.  For every balloon he buys for $n$ dollars, he can buy another for $\frac{2n}{3}$ dollars.  This means it costs him $\frac{5n}{3}$ dollars to buy a bundle of $2$ balloons.  With $30n$ dollars, he can buy $\frac{30n}{\frac{5n}{3}} = 18$ sets of two balloons, so the total number of balloons he can buy is $18\times2  \implies \boxed{\textbf{(C)}\ 36 }$




\textbf{Solution 3}

Similar to solution 1, but quicker. We see the fraction $\frac13,$ so we let each balloon cost $\$3$ WLOG. Thus, Orvin has $\$3\cdot30=\$90.$ Each pair of balloons (one full-priced and one a third off) costs $\$3+\left(1-\frac13\right)\cdot\$3=\$3+\$2=\$5.$ Therefore, Orvin can buy $\frac{\$90}{\$5}=18$ pairs of balloons, at maximum. $18\cdot2=36$ balloons. ~Technodoggo

 


 


	\item (\textit{AMC 12 12B 2014}) Randy drove the first third of his trip on a gravel road, the next $20$ miles on pavement, and the remaining one-fifth on a dirt road. In miles, how long was Randy's trip?


 $\textbf{(A)}\ 30\qquad\textbf{(B)}\ \frac{400}{11}\qquad\textbf{(C)}\ \frac{75}{2}\qquad\textbf{(D)}\ 40\qquad\textbf{(E)}\ \frac{300}{7}$

 \textbf{Answer: }E



\textbf{Solution}

If the first and last legs of his trip account for $\frac{1}{3}$ and $\frac{1}{5}$ of his trip, then the middle leg accounts for $1 - \frac{1}{3} - \frac{1}{5} = \frac{7}{15}$ths of his trip.  This is equal to $20$ miles.  Letting the length of the entire trip equal $x$, we have 
 \[\frac{7}{15}x = 20 \implies x=\boxed{\textbf{(E)}\ \frac{300}{7}}\]



	\item (\textit{AMC 12 12B 2014}) Susie pays for $4$ muffins and $3$ bananas. Calvin spends twice as much paying for $2$ muffins and $16$ bananas. A muffin is how many times as expensive as a banana?


 $\textbf{(A)}\ \frac{3}{2}\qquad\textbf{(B)}\ \frac{5}{3}\qquad\textbf{(C)}\ \frac{7}{4}\qquad\textbf{(D)}\ 2\qquad\textbf{(E)}\ \frac{13}{4}$

 \textbf{Answer: }B



\textbf{Solution}

Let $m$ stand for the cost of a muffin, and let $b$ stand for the value of a banana. We need to find $\frac{m}{b}$, the ratio of the price of the muffins to that of the bananas. We have
 \[2(4m + 3b) = 2m + 16b\]
 \[6m = 10b\]
 \[\frac{m}{b} = \boxed{\textbf{(B)}\ \frac{5}{3}}\]



	\item (\textit{AMC 12 12B 2014}) Doug constructs a square window using $8$ equal-size panes of glass, as shown. The ratio of the height to width for each pane is $5 : 2$, and the borders around and between the panes are $2$ inches wide. In inches, what is the side length of the square window?

                    \begin{figure}[H]
                        \centering
                        \includegraphics[width=0.5\textwidth]{Problem/AMC_12_12B_2014_5_Problem_Figure_1.png}
                    \end{figure}
                    
$\textbf{(A)}\ 26\qquad\textbf{(B)}\ 28\qquad\textbf{(C)}\ 30\qquad\textbf{(D)}\ 32\qquad\textbf{(E)}\ 34$

 \textbf{Answer: }A



\textbf{Solution}

Let the height of the panes equal $5x$, and let the width of the panes equal $2x$.  Now notice that the total width of the borders equals $10$, and the total height of the borders is $6$.  We have
 \[10 + 4(2x) = 6 + 2(5x)\]
 \[x = 2\]Now, the total side length of the window equals 
 \[10+ 4(2x) = 10 + 16 = \boxed{\textbf{(A)}\ 26}\]



	\item (\textit{AMC 12 12B 2014}) Ed and Ann both have lemonade with their lunch. Ed orders the regular size. Ann gets the large lemonade, which is 50\% more than the regular. After both consume $\frac{3}{4}$ of their drinks, Ann gives Ed a third of what she has left, and 2 additional ounces. When they finish their lemonades they realize that they both drank the same amount. How many ounces of lemonade did they drink together?


 $\textbf{(A)}\ 30\qquad\textbf{(B)}\ 32\qquad\textbf{(C)}\ 36\qquad\textbf{(D)}\ 40\qquad\textbf{(E)}\ 50$

 \textbf{Answer: }D



\textbf{Solution}

Let the size of Ed's drink equal $x$ ounces, and let the size of Ann's drink equal $\frac{3}{2}x$ ounces.  After both consume $\frac{3}{4}$ of their drinks, Ed and Ann have $\frac{x}{4}$  and $\frac{3x}{8}$ ounces of their drinks remaining.  Ann gives away $\frac{x}{8} + 2$ ounces to Ed.

 In the end, Ed drank everything in his original lemonade plus what Ann gave him, and Ann drank everything in her original lemonade minus what she gave Ed. Thus we have
 \[x + \frac{x}{8} + 2 = \frac{3x}{2} - \frac{x}{8} - 2\]
 \[x = 16\]The total amount the two of them drank is simply 
 \[x + \frac{3}{2}x = 16 + 24 = \boxed{\textbf{(D)}\ 40}\]



	\item (\textit{AMC 12 12B 2014}) For how many positive integers $n$ is $\frac{n}{30-n}$ also a positive integer?


 $\textbf{(A)}\ 4\qquad\textbf{(B)}\ 5\qquad\textbf{(C)}\ 6\qquad\textbf{(D)}\ 7\qquad\textbf{(E)}\ 8$

 \textbf{Answer: }D



\textbf{Solution 1}

We know that $n \le 30$ or else $30-n$ will be negative, resulting in a negative fraction.  We also know that $n \ge 15$ or else the fraction's denominator will exceed its numerator making the fraction unable to equal a positive integer value.  Substituting all values $n$ from $15$ to $30$ gives us integer values for $n=15, 20, 24, 25, 27, 28, 29$.  Counting them up, we have $\boxed{\textbf{(D)}\ 7}$ possible values for $n$.  

 



\textbf{Solution 2}

Let $m = \frac{n}{30-n}$ be a positive integer ($m \ge 1$). Rearranging gives
 \[n = m(30 - n) \implies mn - 30m + n = 0\]To factor using Simon's Favorite Factoring Trick, add $-30$ to both sides:
 \[mn - 30m + n - 30 = -30 \implies (m + 1)(n - 30) = -30\]Multiplying both sides by $-1$, we get
 \[(m + 1)(30 - n) = 30\]Since $m$ is a positive integer ($m \ge 1$), $m + 1 \ge 2$. Therefore, $m + 1$ can be any positive divisor of $30$ except $1$.

 The prime factorization of $30$ is $2^1 \cdot 3^1 \cdot 5^1$, which has $(1 + 1)(1 + 1)(1 + 1) = 8$ positive divisors. Excluding $m + 1 = 1$ (which corresponds to $n = 0$, violating $n > 0$) leaves $8 - 1 = \mathbf{7}$ valid solutions, so the answer is $\boxed{\textbf{(D)}\ 7}$.

 



\textbf{Solution 3}

Let $\frac{n}{30-n} = m$, where $m\in\mathbb{N}$. Solving for $n$, we find that $n=\frac{30m}{m+1}$. Because $m$ and $m+1$ are relatively prime, $m+1\mid 30$. Our answer is the number of proper divisors of $2^1 3^1 5^1$, which is $(1+1)(1+1)(1+1)-1=\boxed{\textbf{(D)}\ 7}$.

 



\textbf{Solution 4}

First of all, we must have $n\leq{29}$ because the denominator will be $0$ if $n=30$ and the fraction will be negative if $n\geq{31}$.

 
 \[30-n\geq{1}\]

 \[\frac{n}{30-n} = \frac{n-30+30}{30-n} = \frac{n-30}{30-n}+\frac{30}{30-n} = \frac{30}{30-n}-1\]

 Thus, $30-n$ must be a positive divisor of $30$.

 $30=2\times 3\times 5$ has $8$ positive divisors. However, $30-n=30$ is not obtainable because this will make $n$ equal to $0$.

 The seven other divisors of $30$, namely $1,2,3,5,6,10$ and $15$, can all be assigned to the value of $30-n$.

 The answer to the problem is $\boxed{\textbf{(D)}\ 7}$.

 



\textbf{Solution 5}

We know that $30-n\mid n$. Then, by divisibility rules:
 \[30-n\mid n+30-n\]
 \[30-n\mid 30\]

 There are $8$ divisors of $30$, but $n$ must be positive, so $30\mid 30$ isn't counted, meaning we have $\boxed{\textbf{(D)}\ 7}$.

 



\textbf{Solution 6}

We recognize that $15\le n < 30$ because for $\frac{n}{30-n}$ to be a positive integer, $30-n \le n \implies 2n \ge 30 \implies n \ge 15$. Testing the numbers in this range gives $n \in \{15, 20, 24, 25, 27, 28, 29\}$, yielding $\boxed{\textbf{(D)}\ 7}$.

 


	\item (\textit{AMC 12 12B 2014}) In the addition shown below $A$, $B$, $C$, and $D$ are distinct digits. How many different values are possible for $D$?


 
 \[\begin{tabular}{cccccc}&A&B&B&C&B\\ +&B&C&A&D&A\\ \hline &D&B&D&D&D\end{tabular}\]
$\textbf{(A)}\ 2\qquad\textbf{(B)}\ 4\qquad\textbf{(C)}\ 7\qquad\textbf{(D)}\ 8\qquad\textbf{(E)}\ 9$

 \textbf{Answer: }C



\textbf{Solution 1}

From the first column, we see $A+B < 10$ because it yields a single digit answer.  From the fourth column, we see that $C+D$ equals $D$ and therefore $C = 0$.  We know that $A+B = D$.  Therefore, the number of values $D$ can take is equal to the number of possible sums less than $10$ that can be formed by adding two distinct natural numbers.  Letting $A=1$, and letting $B=2,3,4,5,6,7,8$, we have 
 \[D = 3,4,5,6,7,8,9 \implies \boxed{\textbf{(C)}\ 7}\]




\textbf{Solution 2}

First of all, focus on the second column from right to left. If we set $C=0$, $C+D=D$ will be easily satisfied. Hence, we will first try setting $C$ to $0$.

 
The addition has become

 

 \[\begin{tabular}{cccccc}&A&B&B&0&B\\ +&B&0&A&D&A\\ \hline &D&B&D&D&D\end{tabular}\]

 
It is now in a very regular form —— Two columns are automatically satisfied, and the other three columns are all in the form $A+B=D$. As long as we have $A+B=D$, all the five columns will be satisfied.

 
Setting $A=1$ and $B=2,3,4,5,6,7,8$, we can construct any value of $D$ in $3,4,5,6,7,8,9$. But what about $D=0,1$ or $2$?

 
We can actually prove that they are all impossible. We now dive into the fifth column from right to left.

 
If there is no carrying, $A+B=D$. If carrying occurs, $A+B+1=D$. No matter whether carrying takes place or not, we have the inequality $A+B\leq{D}$.

 
If $D=0$, $A=B=0$. This is a contradiction as all the digits are distinct. $D\neq{0}$.

 
If $D=1$, $A+B\leq{1}$. One of $A$ and $B$ has to be zero.

 
\textit{Lemma:}

 
Neither $A$ nor $B$ can be zero.

 
\textit{Proof:}

 
Focus on the first column from right to left.

 If $A=0$, $B=D$. The digits are no longer distinct.

 If $B=0$, $A=D$. The digits are also no longer distinct.

 
Thus, $D\neq{1}$.

 
If $D=2$, $A+B\leq{2}$. But $(A,B)\neq(1,1)$ as the digits are distinct. $(A,B)=(0,2)$ or $(2,0)$, which both involve one of $A$ and $B$ being zero.

 
None of $0,1$ and $2$ can be a value of $D$. This completes the proof.

 
The answer to the problem is $7$.

 



\textbf{Solution 3 (Equation Algorithm)}

It is intuitively obvious, even to the most casual observer that the problem statement can be rewritten as:

 $10^4A + 10^4B + 10^3B + 10^3C + 10^2B + 10^2A + 10C + 10D + B + A = 10^4D + 10^3B + 10^2D + 10D + D$. This equation can be simplified into:

 $10^4A + 10^4B + 10^3C + 10^2B + 10^2A + 10C + B + A = 10^4D + 10^2D + D$.

 Now from here, it should hopefully make sense that $A + B = D$ by looking at the one's digit of both equations. Factoring out $A + B$ gives:

 $10^4(A+B) + 10^3C + 10^2(A+B) + 10C + (B + A) = 10^4D + 10^2D + D$. 

 Which equals:$10^4(D) + 10^3C + 10^2(D) + 10C + D = 10^4D + 10^2D + D$. 

 This simplifies into:$10^3C + 10C = 0$.

 Therefore $C = 0$.

 This means that $A + B = D$ and $D < 10$ or else there would be parts carried over in the equation. The positive integers that satisfy this equation are a minimum $(2, 1)$ and a maximum of $(4, 5)$. This means that $D = 3$ $, 4$ $, 5$ $, 6$ $, 7$ $, 8$ $, 9$. Giving $\boxed{\textbf{(C)}\ 7}$



	\item (\textit{AMC 12 12B 2014}) Convex quadrilateral $ABCD$ has $AB=3$, $BC=4$, $CD=13$, $AD=12$, and $\angle ABC=90^{\circ}$, as shown. What is the area of the quadrilateral?


 
                    \begin{figure}[H]
                        \centering
                        \includegraphics[width=0.5\textwidth]{Problem/AMC_12_12B_2014_9_Problem_Figure_1.png}
                    \end{figure}
                    
$\textbf{(A)}\ 30\qquad\textbf{(B)}\ 36\qquad\textbf{(C)}\ 40\qquad\textbf{(D)}\ 48\qquad\textbf{(E)}\ 58.5$

 \textbf{Answer: }B



\textbf{Solution}

Note that by the pythagorean theorem, $AC=5$.  Also note that $\angle CAD$ is a right angle because $\triangle CAD$ is a right triangle.  The area of the quadrilateral is the sum of the areas of $\triangle ABC$ and $\triangle CAD$ which is equal to 
 \[\frac{3\times4}{2} + \frac{5\times12}{2} = 6 + 30 = \boxed{\textbf{(B)}\ 36}\]



	\item (\textit{AMC 12 12B 2014}) Danica drove her new car on a trip for a positive whole number of hours, averaging 55 miles per hour. At the beginning of the trip, $abc$ miles was displayed on the odometer, where $abc$ is a 3-digit number with $a \geq{1}$ and $a+b+c \leq{7}$. At the end of the trip, the odometer showed $cba$ miles. What is $a^2+b^2+c^2?$.


 $\textbf{(A)}\ 26\qquad\textbf{(B)}\ 27\qquad\textbf{(C)}\ 36\qquad\textbf{(D)}\ 37\qquad\textbf{(E)}\ 41$

 \textbf{Answer: }D



\textbf{Solution 1}

We know that the number of miles she drove is divisible by $5$, so $a$ and $c$ must either be equal or differ by $5$. We can quickly conclude that the former is impossible; otherwise, she would have travelled 0 miles. Therefore, $a$ and $c$ must be $5$ apart.  Because we know that $c > a$ and $a + c \le 7$ and $a \ge 1$, we find that the only possible values for $a$ and $c$ are $1$ and $6$, respectively.  Because $a + b + c \le 7$, $b = 0$.  Therefore, we have
 \[a^2 + b^2 + c^2 = 36 + 0 + 1 = \boxed{\textbf{(D)}\ 37}\]




\textbf{Solution 2}

Let the number of hours Danica drove be $k$. Then we know that $100a + 10b + c + 55k$ = $100c + 10b + a$. Simplifying, we have $99c - 99a = 55k$, or $9c - 9a = 5k$. Thus, k is divisible by $9$. Because $55 * 18 = 990$, $k$ must be $9$, and therefore $c - a = 5$. Because $a + b + c \leq{7}$ and $a \geq{1}$, $a = 1$, $c = 6$ and $b = 0$, and our answer is $a^2 + b^2 + c^2 = 6^2 + 0^2 + 1^2 = 37$, or $\boxed{D}$.

 


 


	\item (\textit{AMC 12 12B 2014}) A list of 11 positive integers has a mean of 10, a median of 9, and a unique mode of 8. What is the largest possible value of an integer in the list?


 $\textbf{(A)}\ 24\qquad\textbf{(B)}\ 30\qquad\textbf{(C)}\ 31\qquad\textbf{(D)}\ 33\qquad\textbf{(E)}\ 35$

 \textbf{Answer: }E



\textbf{Solution}

We start off with the fact that the median is $9$, so we must have $a, b, c, d, e, 9, f, g, h, i, j$, listed in ascending order. Note that the integers do not have to be distinct. 

 Since the mode is $8$, we have to have at least $2$ occurrences of $8$ in the list. If there are $2$ occurrences of $8$ in the list, we will have $a, b, c, 8, 8, 9, f, g, h, i, j$. In this case, since $8$ is the unique mode, the rest of the integers have to be distinct. So we minimize $a,b,c,f,g,h,i$ in order to maximize $j$. If we let the list be $1,2,3,8,8,9,10,11,12,13,j$, then $j = 11 \times 10 - (1+2+3+8+8+9+10+11+12+13) = 33$. 

 Next, consider the case where there are $3$ occurrences of $8$ in the list. Now, we can have two occurrences of another integer in the list. We try $1,1,8,8,8,9,9,10,10,11,j$. Following the same process as above, we get $j = 11 \times 10 - (1+1+8+8+8+9+9+10+10+11) = 35$. As this is the highest choice in the list, we know this is our answer. Therefore, the answer is $\boxed{\textbf{(E) }35}$



	\item (\textit{AMC 12 12B 2014}) A set $S$ consists of triangles whose sides have integer lengths less than 5, and no two elements of $S$ are congruent or similar. What is the largest number of elements that $S$ can have?


 $\textbf{(A)}\ 8\qquad\textbf{(B)}\ 9\qquad\textbf{(C)}\ 10\qquad\textbf{(D)}\ 11\qquad\textbf{(E)}\ 12$

 \textbf{Answer: }B



\textbf{Solution 1}

Define $T$ to be the set of all integral triples $(a, b, c)$ such that $a \ge b \ge c$, $b+c > a$, and $a, b, c < 5$. Now we enumerate the elements of $T$:

 $(4, 4, 4)$
$(4, 4, 3)$
$(4, 4, 2)$
$(4, 4, 1)$
$(4, 3, 3)$
$(4, 3, 2)$
$(3, 3, 3)$
$(3, 3, 2)$
$(3, 3, 1)$
$(3, 2, 2)$
$(2, 2, 2)$
$(2, 2, 1)$
$(1, 1, 1)$

 It should be clear that $|S|$ is simply $|T|$ minus the larger "duplicates" (e.g. $(2, 2, 2)$ is a larger duplicate of $(1, 1, 1)$). Since $|T|$ is $13$ and the number of higher duplicates is $4$, the answer is $13 - 4$ or $\boxed{\textbf{(B)}\ 9}$.

 


 



\textbf{Solution 2}

First of all, we can have at most one equilateral triangle in our set because all equilateral triangles are similar. We may pick the one with side lengths $(1,1,1)$.

 
The isosceles non-equilateral triangles are composed of:

 $(1,2,2),(1,3,3),(1,4,4),(2,3,3),(2,4,4),(3,2,2),(3,4,4),(4,3,3)$ (ordered by base length)

 
The triplets $(2,1,1),(3,1,1),(4,1,1),(4,2,2)$ are not capable of forming triangles.

 
We can include at most $7$ isosceles non-equilateral triangles in our set. If we include all $8$ triangles, $(1,2,2)$ and $(2,4,4)$ will be similar.

 
Finally, there exists only one scalene triangle —— $(2,3,4)$. This is the only such triangle because a triangle cannot be formed if we use $1$ as the shortest side.

 
A triangle of one type will never clash with a triangle of another type under this categorization.

 
The answer to the problem is $1+7+1=9$.

 


	\item (\textit{AMC 12 12B 2014}) Let $b$ be a positive real number greater than $1$. Find the least possible value of $b$ such that: There exists a real number $a\in(1,b)$ for which there exists neither a triangle of side lengths $1$, $a$ and $b$ with positive area, nor a triangle of side lengths $\frac{1}{b}$, $\frac{1}{a}$ and $1$ with positive area.


 $\textbf{(A)}\ \frac{3+\sqrt{3}}{2}\qquad\textbf{(B)}\ \frac{5}{2}\qquad\textbf{(C)}\ \frac{3+\sqrt{5}}{2}\qquad\textbf{(D)}\ \frac{3+\sqrt{6}}{2}\qquad\textbf{(E)}\ 3$

 \textbf{Answer: }C



\textbf{Solution 1}

Notice that $1>\frac{1}{a}>\frac{1}{b}$.  Using the triangle inequality, we find 
 \[a+1 > b \implies a>b-1\]
 \[\frac{1}{a}+\frac{1}{b} > 1\]In order for us the find the lowest possible value for $b$, we try to create two degenerate triangles where the sum of the smallest two sides equals the largest side.Thus we get 
 \[a=b-1\] and 
 \[\frac{1}{a} + \frac{1}{b}=1\]Substituting, we get
 \[\frac{1}{b-1}+\frac{1}{b}=\frac{b+b-1}{b(b-1)}=1\]
 \[\frac{2b-1}{b(b-1)} = 1\]
 \[2b-1=b^2-b\]Solving for $b$ using the quadratic equation, we get
 \[b^2-3b+1=0 \implies b = \boxed{\textbf{(C)} \ \frac{3+\sqrt{5}}{2}}\]




\textbf{Solution 2 (similar to Solution 1)}

We can form degenerate triangles from the given information. Since $b>a>1$, we know that $a,b,1$ must satisfy $a + 1 \le b$ to be degenerate in the first scenario. 

 
Similarly, since $\frac{1}{n} + 1 > 1$ for all positive integers $n$, the only condition that is guaranteed to form a degenerate triangle from the second instance is $\frac{1}{a} + \frac{1}{b} \le 1$.

 
Solving the second inequality for $a$ in terms of $b$ (for more-easily apparent inequality manipulation) yields $b\ge \frac{a}{a-1}$. Notice that $b> a+1 \ge \frac{a}{a-1}$ for almost all positive real numbers $a>1$ (see note at the end). 

 
Thus, it suffices to solve the inequality $a+1 \ge \frac{a}{a-1}$ for $a$ whose only positive solution is $a \ge \frac{1 + \sqrt{5}}{2}$. 

 
Substituting back into the first inequality yields $\boxed{b \ge \frac{3 + \sqrt{5}}{2}}$ and thus the answer is $\text{C}$.

 


 Note: this solution assumes that $a$ is not less than $1.618$, the points at which the last inequality. 

 


 



\textbf{Solution 3}

By the triangle inequality we have $a+b \ge 1$, $a + 1 \ge b$, $b+1 \ge a$ and $b(1+a) \ge a$, $a(1+b) \ge b$, $a+b \ge ab$. Clearly a triangle can't have negative area so, our triangle must have $0$ area which means equality must be obtained for all the inequalities above. Since $a = 1+b$ we know that $(1+b)^2 = b$ which simplifies to $b^2 - 3b + 1 = 0$. Solving for $b$ we get $\frac{3 + \sqrt{5}}{2}$. So the answer is C.~coolmath\_2018

 



\textbf{Solution 4}

$1+a \leq b, \quad \frac{1}{b} + \frac{1}{a} \leq 1$
$\implies a \leq b-1, \quad \frac{a+b}{ab} \leq 1$
$\implies a+b \leq ab \implies b \leq ab-a \implies b \leq(b-1)a \implies a \geq \frac{b}{b-1}$
$\therefore \frac{b}{b-1} \leq a \leq b-1$
$\frac{b}{b-1} \leq b-1 \implies b \leq b^2-2b+1 \implies b^2-3b+1 \geq 0$
$b= \frac{3 \pm \sqrt{9-4}}{2}= \frac{3 \pm \sqrt{5}}{2}$
$b=\frac{3 + \sqrt{5}}{2}$

 select C

 


 



\textbf{Solution 5 (Rigorous)}

$1<a<b$
$\frac{1}{b}<\frac{1}{a}<1$

 
Three sides cannot form a triangle if and only if the sum of the lengths of the two shorter sides does not exceed the length of the longest side.

 It is given that

 
 \[1+a\leq{b}\]
 \[\frac{1}{b}+\frac{1}{a}\leq{1}\]

 The second inequality is equivalent to $a+b\leq{ab}$
$a-ab\leq{-b}$
$a(1-b)\leq{-b}$
$a\geq\frac{-b}{1-b}$ (since $1-b$ is negative)

 $a\geq\frac{b}{b-1}$

 
Substituting $a\geq\frac{b}{b-1}$ into the first inequality, we obtain $b\geq{1+a}\geq{1+\frac{b}{b-1}}$.

 
$b\geq{1+\frac{b}{b-1}}$
$b\geq{\frac{b-1}{b-1}+\frac{b}{b-1}}$
$b\geq{\frac{2b-1}{b-1}}$
$b(b-1)\geq{2b-1}$ (since $b-1$ is positive)

 $b^2-b\geq{2b-1}$
$b^2-3b+1\geq{0}$
$b$ has to be greater or equal to both of the solutions to this quadratic equation or less than or equal to both of them. We have either $b\geq{\frac{3+{\sqrt{5}}}{2}}$ or $b\leq{\frac{3-{\sqrt{5}}}{2}}$. The latter is impossible because it will make $b$ less than $1$.

 
$b\geq{\frac{3+{\sqrt{5}}}{2}}$ is a necessary condition. We have found a lower bound for $b$. Now we need to prove that $b={\frac{3+{\sqrt{5}}}{2}}$ can satisfy the conditions.

 
A value of $a$ that is just adequate to make a triangle degenerate might be useful. We choose $a=b-1={\frac{1+{\sqrt{5}}}{2}}$.

 
We have $1+\frac{1+{\sqrt{5}}}{2}=\frac{3+{\sqrt{5}}}{2}$, which means that there is not a triangle with positive area with these side lengths.

 We also have $1\div\frac{3+{\sqrt{5}}}{2}+1\div\frac{1+{\sqrt{5}}}{2}$
$=\frac{2}{3+{\sqrt{5}}}+\frac{2}{1+{\sqrt{5}}}$
$=\frac{2({\sqrt{5}}-3)}{({\sqrt{5}}+3)({\sqrt{5}}-3)}+\frac{2({\sqrt{5}}-1)}{({\sqrt{5}}+1)({\sqrt{5}}-1)}$
$=\frac{2({\sqrt{5}}-3)}{-4}+\frac{2({\sqrt{5}}-1)}{4}$
$=\frac{3-\sqrt{5}}{2}+\frac{{\sqrt{5}}-1}{2}$
$=1$, which means that there is not a triangle with positive area with these side lengths.

 
We have shown that $b$ has a lower bound of ${\frac{3+{\sqrt{5}}}{2}}$, and for $b={\frac{3+{\sqrt{5}}}{2}}$, there exists a value of $a$ that satisfies the conditions. This completes the proof.

 
The smallest possible value of $b$ is ${\frac{3+{\sqrt{5}}}{2}}$.

 


	\item (\textit{AMC 12 12B 2014}) A rectangular box has a total surface area of 94 square inches. The sum of the lengths of all its edges is 48 inches. What is the sum of the lengths in inches of all of its interior diagonals?


 $\textbf{(A)}\ 8\sqrt{3}\qquad\textbf{(B)}\ 10\sqrt{2}\qquad\textbf{(C)}\ 16\sqrt{3}\qquad\textbf{(D)}\ 20\sqrt{2}\qquad\textbf{(E)}\ 40\sqrt{2}$

 \textbf{Answer: }D



\textbf{Solution 1}

Let the side lengths of the rectangular box be $x, y$  and $z$. From the information we get 

 
 \[4(x+y+z) = 48 \Rightarrow x+y+z = 12\]

 \[2(xy+yz+xz) = 94\]

 The sum of all the  lengths of the box's interior diagonals is 

 
 \[4 \sqrt{x^2+y^2+z^2}\]

 Squaring the first expression, we get: 

 \begin{align*}144 =(x+y+z)^2 &=  x^2+y^2+z^2 + 2(xy+yz+xz) \\&=  x^2+y^2+z^2 + 94.\end{align*}


 Hence 
 \[x^2+y^2+z^2 = 50\]

 \[4 \sqrt{x^2+y^2+z^2} =  \boxed{\textbf{(D)}\ 20\sqrt 2}\]




\textbf{Solution 2 (Observation)}

Observe that a rectangular prism with side lengths of $3,4,5$ fits all the necessary conditions. Thus, the answer is 
 \[4\sqrt{3^2+4^2+5^2} = \boxed{\textbf{(D)}\ 20\sqrt 2}\]~LuisFonseca123

 



\textbf{Solution 3 (Estimation?)}

Observe that a cube with side length $4$ has an equal sum of edge lengths and a similar surface area $(96)$. We can evaluate the total interior diagonal length of this cube to be $16\sqrt{3}$. Since the cube almost follows the problem's conditions, we can infer that the total interior diagonal length of the rectangular prism to be close but unequal to $16\sqrt{3}$, which leads us to an answer of 
 \[\boxed{\textbf{(D)}\ 20\sqrt 2}\]




\textbf{Solution 4 (Guessinggggggg)}

Observe that it just cannot be E, because this is way out of the others range. Next, observe that it cannot be B because it is not only pretty far away from C, but it is also the only one that is not divisible by 4. So we have 10, 16, 20 left. Note that the average of all the answers is highly skewed to the right of the median, therefore the answer will be around there (this is because the average of the answers is a good answer for MAA to use because it means that it is not obviously lower or higher than all 4 of the other answers, plus it is easier to get mislead with B or D.) With that being said, we can conclude that the answer is almost certainly D, because the average will be somewhat close to 
 \[\boxed{\textbf{(D)}\ 20\sqrt 2}\]



	\item (\textit{AMC 12 12B 2014}) When $p = \sum\limits_{k=1}^{6} k \ln{k}$, the number $e^p$ is an integer.  What is the largest power of $2$ that is a factor of $e^p$ ?


 $\textbf{(A)}\ 2^{12}\qquad\textbf{(B)}\ 2^{14}\qquad\textbf{(C)}\ 2^{16}\qquad\textbf{(D)}\ 2^{18}\qquad\textbf{(E)}\ 2^{20}$

 \textbf{Answer: }C



\textbf{Solution}

Let's write out the sum.  Our sum is equal to 
 \[1 \ln{1} + 2 \ln{2} + 3 \ln{3} + 4 \ln{4} + 5 \ln{5} + 6 \ln{6} =\]
 \[\ln{1^1} + \ln{2^2} + \ln{3^3} + \ln {4^4} + \ln{5^5} + \ln {6^6} =\]
 \[\ln{(1^1\times2^2\times3^3\times4^4\times5^5\times6^6)}\]Raising $e$ to the power of this quantity eliminates the natural logarithm, which leaves us with 
 \[e^p = 1^1\times2^2\times3^3\times4^4\times5^5\times6^6\]This product has $2$ powers of $2$ in the $2^2$ factor, $4*2=8$ powers of $2$ in the $4^4$ factor, and $6$ powers of $2$ in the $6^6$ factor.This adds up to $2+8+6=16$ powers of two which divide into our quantity, so our answer is $\boxed{\textbf{(C)}\ 2^{16}}$




\textbf{Solution 2 (less log knowledge)}

You could use a different method of solving this, without using specialized log rules but you do have to know that the summation of logs equals the products. Using this, we can see that 1ln1 + 2ln2 + 3ln3 + 4ln4 + 5ln5 + 6ln6 = (ln1+ln2+ln3+ln4+ln5+ln6) + (ln2+ln3...ln6) + (ln3...ln6) + (ln4...ln6) + (ln5+ln6) + (ln6)So we see it equates to ln(6!)+ln(6!)+ln(6!/2!)+ln(6!/3!)+ln(6!/4!)+ln(6!/5!), or ln(6!\textasciicircum{}6/(2*3!*4!*5!). If we raise e to this, we eliminate the ln part to get solely the argument. We now see how many powers of two 6! or 720 contains, which is 2\textasciicircum{}4. So we have on the numerator 24 powers of two and on the denominator 8 powers of two. This gets us 16 powers of two, which is answer choice $\boxed{\textbf{(C)}\ 2^{16}}$

 -kingoringo828 (please fix latex)

 


	\item (\textit{AMC 12 12B 2014}) Let $P$ be a cubic polynomial with $P(0) = k$, $P(1) = 2k$, and $P(-1) = 3k$.  What is $P(2) + P(-2)$ ?


 $\textbf{(A)}\ 0\qquad\textbf{(B)}\ k\qquad\textbf{(C)}\ 6k\qquad\textbf{(D)}\ 7k\qquad\textbf{(E)}\ 14k$

 \textbf{Answer: }E



\textbf{Solution}

Let $P(x) = Ax^3+Bx^2 + Cx+D$.   Plugging in $0$ for $x$, we find $D=k$, and plugging in $1$ and $-1$ for $x$, we obtain the following equations:
 \[A+B+C+k=2k\]
 \[-A+B-C+k=3k\]Adding these two equations together, we get
 \[2B=3k\]If we plug in $2$ and $-2$ in for $x$, we find that 
 \[P(2)+P(-2) = 8A+4B+2C+k+(-8A+4B-2C+k)=8B+2k\]Multiplying the third equation by $4$ and adding $2k$ gives us our desired result, so 
 \[P(2)+P(-2)=12k+2k=\boxed{\textbf{(E)}\ 14k}\]

 A construction that satisfies all the conditions can be $P(x)=-x^3+3x^2+2$ and $k=2$.

 $P(0)=2=k$
$P(1)=-1+3+2=4=2k$
$P(-1)=1+3+2=6=3k$
$P(2)=-8+12+2=6=3k$
$P(-2)=8+12+2=22=11k$
$P(2)+P(-2)=14k$




\textbf{Solution 2}

If we use \href{https://artofproblemsolving.com/wiki/index.php?title=Gregory\%27s\_Triangle}{Gregory's Triangle}, the following happens:
 \[P(-1), P(0), P(1)\]
 \[3k ,  k  ,  2k\]
 \[-2k , k\]
 \[3k\]

 Since this is cubic, the common difference is $3k$ for the linear level so the string of $3k$s are infinite in each direction.If we put a $3k$ on each side of the original $3k$, we can solve for $P(-2)$ and $P(2)$.

 
 \[P(-2), P(-1), P(0), P(1), P(2)\]
 \[8k ,   3k ,  k  ,  2k ,  6k\]
 \[-5k  , -2k ,  k  ,  4k\]
 \[3k   ,  3k ,  3k\]

 The above shows us that $P(-2)$ is $8k$ and $P(2)$ is $6k$ so $8k+6k=14k$.

 NOTE (not from author): The link you put for gregory's triangle doesn't work so please explain it in your post or find a resource that does work; there isn't much on google.

 NOTE (not from author or user above): I have now updated the link. It should work.

 



\textbf{Solution 3}

First, define 
 \[G(x)=P(x)+P(-x).\] We have 
 \[G(-1)=G(1)=5k\] and 
 \[G(0)=2k.\] Notice that 
 \[G(x)\] is of the form 
 \[a*x^2+b,\] since if we added 
 \[P(x)+P(-x),\] the 
 \[x\] and 
 \[x^3\] terms would cancel out. Plug in the values, and you get 
 \[a=3k, b=2k,\] so 
 \[P(2)+P(-2)=G(2)=14k.\]

 (E)

 


	\item (\textit{AMC 12 12B 2014}) Let $P$ be the parabola with equation $y=x^2$ and let $Q = (20, 14)$. There are real numbers $r$ and $s$ such that the line through $Q$ with slope $m$ does not intersect $P$ if and only if $r < m < s$. What is $r + s$?


 $\textbf{(A)}\ 1\qquad\textbf{(B)}\ 26\qquad\textbf{(C)}\ 40\qquad\textbf{(D)}\ 52\qquad\textbf{(E)}\ 80$

 \textbf{Answer: }E



\textbf{Solution 1 (Algebra Based)}

Let $y = m(x - 20) + 14$. Equating them:

 $x^2 = mx - 20m + 14$
$x^2 - mx + 20m - 14 = 0$

 For there to be no solutions, the discriminant must be less than zero:

 $m^2 - 4(20m - 14) < 0$
$m^2 - 80m + 56 < 0$.

 So $m < 0$ for $r < m < s$ where $r$ and $s$ are the roots of $m^2 - 80m + 56 = 0$ and their sum by Vieta's formulas is $\boxed{\textbf{(E)}\ 80}$.

 



\textbf{Solution 2 (Calculus-based)}

The line will begin to intercept the parabola when its slope equals that of the parabola at the point of tangency.  Taking the derivative of the equation of the parabola, we get that the slope equals $2x$.  Using the slope formula, we find that the slope of the tangent line to the parabola also equals $\frac{14-x^2}{20-x}$.  Setting these two equal to each other, we get
 \[2x = \frac{14-x^2}{20-x} \implies x^2-40x+14 = 0\]Solving for $x$, we get 
 \[x= 20\pm \sqrt{386}\]The sum of the two possible values for $x$ where the line is tangent to the parabola is $40$, and the sum of the slopes of these two tangent lines is equal to $2x$, or $\boxed{\textbf{(E)}\ 80}$.

 



\textbf{Solution 3 (fake)}

The smaller solution is basically negligible in comparison with the solution with the larger slope. Try some values like of $x^2 = y$ where $x = 40$ => $y = 1600$, and slope ~80. Trying a few values leads us to conclude the least possible value is around $80$, so the answer is $\boxed {\textbf{(E)}\ 80}$. 

 


	\item (\textit{AMC 12 12B 2014}) The numbers $1$, $2$, $3$, $4$, $5$, are to be arranged in a circle.  An arrangement is $\textit{bad}$ if it is not true that for every $n$ from $1$ to $15$ one can find a subset of the numbers that appear consecutively on the circle that sum to $n$.  Arrangements that differ only by a rotation or a reflection are considered the same.  How many different bad arrangements are there?


 $\textbf{(A) }1\qquad\textbf{(B) }2\qquad\textbf{(C) }3\qquad\textbf{(D) }4\qquad\textbf{(E) }5$

 \textbf{Answer: }B



\textbf{Solution 1}

We see that there are $5!$ total ways to arrange the numbers. However, we can always rotate these numbers so that, for example, the number 1 is always at the top of the circle. Thus, there are only $4!$ ways under rotation, which is not difficult to list out. We systematically list out all $24$ cases. 

 Now, we must examine if they satisfy the conditions. We can see that by choosing one number at a time, we can always obtain subsets with sums 1, 2, 3, 4, and 5. By choosing the full circle, we can obtain 15. By choosing everything except for 1, 2, 3, 4, and 5, we can obtain subsets with sums of 10, 11, 12, 13, and 14. 

 This means that we now only need to check for 6, 7, 8, and 9. However, once we have found a set summing to 6, we can choose everything else and obtain a set summing to 9, and similarly for 7 and 8. Thus, we only need to check each case for whether or not we can obtain 6 or 7. 

 We find that there are only 4 arrangements that satisfy these conditions. However, each of these is a reflection of another. We divide by 2 for these reflections to obtain a final answer of $\boxed{\textbf {(B) }2}$.

 


	\item (\textit{AMC 12 12B 2014}) A sphere is inscribed in a truncated right circular cone as shown. The volume of the truncated cone is twice that of the sphere. What is the ratio of the radius of the bottom base of the truncated cone to the radius of the top base of the truncated cone?

                    \begin{figure}[H]
                        \centering
                        \includegraphics[width=0.5\textwidth]{Problem/AMC_12_12B_2014_19_Problem_Figure_1.png}
                    \end{figure}
                    
$\text{(A) } \dfrac32 \quad \text{(B) } \dfrac{1+\sqrt5}2 \quad \text{(C) } \sqrt3 \quad \text{(D) } 2 \quad \text{(E) } \dfrac{3+\sqrt5}2$

 \textbf{Answer: }E



\textbf{Solution 1}

First, we draw the vertical cross-section passing through the middle of the frustum.let the top base equal 2 and the bottom base to be equal to 2r
                    \begin{figure}[H]
                        \centering
                        \includegraphics[width=0.5\textwidth]{Solution/AMC_12_12B_2014_19_Solution_Figure_1.png}
                    \end{figure}
                    

 then using the Pythagorean theorem we have:$(r+1)^2=(2s)^2+(r-1)^2$which is equivalent to:$r^2+2r+1=4s^2+r^2-2r+1$subtracting $r^2-2r+1$ from both sides$4r=4s^2$solving for s we get:
 \[s=\sqrt{r}\]next we can find the volume of the frustum and of the sphere and we know $V_{\text{frustum}}=2V_{\text{sphere}}$ so we can solve for $s$using $V_{\text{frustum}}=\frac{\pi*h}{3}(R^2+r^2+Rr)$we get:
 \[V_{\text{frustum}}=\frac{\pi*2\sqrt{r}}{3}(r^2+r+1)\]using  $V_{\text{sphere}}=\dfrac{4s^{3}\pi}{3}$we get 
 \[V_{\text{sphere}}=\dfrac{4(\sqrt{r})^{3}\pi}{3}\]so we have:
 \[\frac{\pi*2\sqrt{r}}{3}(r^2+r+1)=2*\dfrac{4(\sqrt{r})^{3}\pi}{3}\]dividing by $\frac{2\pi*\sqrt{r}}{3}$we get 
 \[r^2+r+1=4r\]which is equivalent to 
 \[r^2-3r+1=0\]$r=\dfrac{3\pm\sqrt{(-3)^2-4*1*1}}{2*1}$so 
 \[r=\dfrac{3+\sqrt{5}}{2}\longrightarrow \boxed{E}\]




\textbf{Solution 2(ADD DIAGRAM)}

Let's once again look at the cross section of the frustum. Let the angle from the center of the sphere to a point on the circumference of the bottom circle be $\theta.$ This implies that the angle from the center of the sphere to a point on the circumference of the top circle is $90 - \theta.$ Hence the bottom radius is $r\tan{\theta}$ and the top radius is $\frac {r}{\tan {\theta}}.$ This means that the radio between the bottom radius and top radius is $(\tan {\theta})^2.$ Using the frustum volume formula, we find that the are of this figure is $\frac{2\pi r}{3}(r^2(\tan {\theta})^2 + r^2 + \frac {r^2} {(\tan {\theta})^2}).$ We can equate this to $\frac {8\pi*r^3} 3.$ Simplifying, we are left with a quadratic conveniently in $(\tan {\theta})^2.$ The quadratic is $(\tan {\theta})^4 - 3(\tan {\theta})^2 + 1 = 0.$ This gives us $(\tan {\theta})^2 = \dfrac{3+\sqrt{5}}{2}\longrightarrow \boxed{E}$

 \textbf{Remark} For problems that involve a circle inscribed into an isosceles trapezoid, the following fact is very useful. If we let the bases be $x$ and $y$, and the height be $h$, then $h = \sqrt{xy}$. Then, we could solve from solution 1 directly knowing the radius in terms of the base. By letting the upper base be $1$ and the lower base be $x$, we can find $r$ in terms of $x$, and solve like in solution one.~Puck\_0

 


	\item (\textit{AMC 12 12B 2014}) For how many positive integers $x$ is $\log_{10}(x-40) + \log_{10}(60-x) < 2$ ?


 $\textbf{(A) }10\qquad \textbf{(B) }18\qquad \textbf{(C) }19\qquad \textbf{(D) }20\qquad \textbf{(E) }\text{infinitely many}\qquad$

 \textbf{Answer: }B



\textbf{Solution 1}

The domain of the LHS implies that 
 \[40<x<60\] Since $\log(a)$ is defined only when $a>0$.
 \[\]Begin from the left-hand side
 \[\log_{10}[(x-40)(60-x)]<2\]

 \[(x-40)(60-x)<100\]
 \[-x^2+100x-2500<0\]
 \[(x-50)^2>0\]
 \[x \not = 50\]Hence, we have integers from 41 to 49 and 51 to 59. There are $\boxed{\textbf{(B)} 18}$ integers.

 



\textbf{Solution 2}

This solution is similar to the first solution, but perhaps more clear. 

 By the properties of logarithms, $\log_{10}(x-40)+\log_{10}(60-x)=\log_{10}((x-40)(60-x))\implies\log_{10}((x-40)(60-x)).$ Therefore, $(x-40)(60-x)<100.$ We can see that this function is concave down (that is, it looks like an upside-down U shape); that is, it has a maximum value. This maximum value is attained at $x=50$, in which $(x-40)(60-x)=100.$ This is the one point on the function that does not satisfy the given condition. 

 We also know that, as noted in solution 1, $x$ must be greater than $40$ and less than $60$, since otherwise, $\log_{10}(x-40)$ or $\log_{10}(60-x)$ would be undefined. Thus, the possible values of $x$ are the positive integers from $41$ to $59,$ inclusive, excluding $50.$ Thus, there are $19-1=18$ total integers. 

 


	\item (\textit{AMC 12 12B 2014}) In the figure, $ABCD$ is a square of side length $1$. The rectangles $JKHG$ and $EBCF$ are congruent. What is $BE$?

                    \begin{figure}[H]
                        \centering
                        \includegraphics[width=0.5\textwidth]{Problem/AMC_12_12B_2014_21_Problem_Figure_1.png}
                    \end{figure}
                    
$\textbf{(A) }\frac{1}{2}(\sqrt{6}-2)\qquad\textbf{(B) }\frac{1}{4}\qquad\textbf{(C) }2-\sqrt{3}\qquad\textbf{(D) }\frac{\sqrt{3}}{6}\qquad\textbf{(E) } 1-\frac{\sqrt{2}}{2}$

 \textbf{Answer: }C



\textbf{Solution 1}

Draw the altitude from $H$ to $AB$ and call the foot $L$. Then $HL=1$. Consider $HJ$. It is the hypotenuse of both right triangles $\triangle HGJ$ and $\triangle HLJ$, and we know $JG=HL=1$, so we must have $\triangle HGJ\cong\triangle JLH$ by Hypotenuse-Leg congruence. From this congruence we have $LJ=HG=BE$.

 Notice that all four triangles in this picture are similar. Also, we have $LA=HD=EJ$. So set $x=LJ=HG=BE$ and $y=LA=HD=EJ$. Now $BE+EJ+JL+LA=2(x+y)=1$. This means $x+y=\frac{1}{2}=BE+EJ=BJ$, so $J$ is the midpoint of $AB$. So $\triangle AJG$, along with all other similar triangles in the picture, is a 30-60-90 triangle, and we have $AG=\sqrt{3} \cdot AJ=\sqrt{3}/2$ and subsequently $GD=\frac{2-\sqrt{3}}{2}=KE$. This means $EJ=\sqrt{3} \cdot KE=\frac{2\sqrt{3}-3}{2}$, which gives $BE=\frac{1}{2}-EJ=\frac{4-2\sqrt{3}}{2}=2-\sqrt{3}$, so the answer is $\textbf{(C)}$.

 



\textbf{Solution 2}

Let $BE = x$.  Let $JA = y$.  Because $\angle FKH = \angle EJK = \angle AGJ = \angle DHG$ and $\angle FHK = \angle EKJ = \angle AJG = \angle DGH$, $\triangle KEJ, \triangle JAG, \triangle GDH, \triangle HFK$ are all similar.  Using proportions and the pythagorean theorem, we find 
 \[EK = xy\]
 \[FK = \sqrt{1-y^2}\]
 \[EJ = x\sqrt{1-y^2}\]Because we know that $BE+EJ+AJ = EK + FK = 1$, we can set up a systems of equations and solving for $x$, we get
 \[x + x\sqrt{1-y^2} + y = 1 \implies x= \frac{1-y}{1+\sqrt{1-y^2}}\]
 \[xy + \sqrt{1-y^2} = 1 \implies x= \frac{1-\sqrt{1-y^2}}{y}\]Now solving for $y$, we get 
 \[\frac{1-y}{1+\sqrt{1-y^2}}=\frac{1-\sqrt{1-y^2}}{y} \implies y(1-y)=(1-\sqrt{1-y^2})(1+\sqrt{1-y^2}) \implies y-y^2=y^2 \implies y=\frac{1}{2}\]Plugging into the second equations with $x$, we get
 \[x= 2\left(1-\sqrt{1-\frac{1}{4}}\right) = 2\left(\frac{2-\sqrt{3}}{2} \right) = \boxed{\textbf{(C)}\ 2-\sqrt{3}}\]




\textbf{Solution 3}

Let $BE = x$, $EK = a$, and $EJ = b$. Then $x^2 = a^2 + b^2$ and because $\triangle KEJ \cong \triangle GDH$ and $\triangle KEJ \sim \triangle JAG$, $\frac{GA}{1} = 1 - a = \frac{b}{x}$. Furthermore, the area of the four triangles and the two rectangles sums to 1:

 
 \[1 = 2x + GA\cdot JA + ab\]

 \[1 = 2x + (1 - a)(1 - (x + b)) + ab\]

 \[1 = 2x + \frac{b}{x}(1 - x - b) + \left(1 - \frac{b}{x}\right)b\]

 \[1 = 2x + \frac{b}{x} - b - \frac{b^2}{x} + b - \frac{b^2}{x}\]

 \[x = 2x^2 + b - 2b^2\]

 \[x - b = 2(x - b)(x + b)\]

 \[x + b = \frac{1}{2}\]

 \[b = \frac{1}{2} - x\]

 \[a = 1 - \frac{b}{x} = 2 - \frac{1}{2x}\]

 By the Pythagorean theorem: $x^2 = a^2 + b^2$

 \[x^2 = \left(2 - \frac{1}{2x}\right)^2 + \left(\frac{1}{2} - x\right)^2\]

 \[x^2 = 4 - \frac{2}{x} + \frac{1}{4x^2} + \frac{1}{4} - x + x^2\]

 \[0 =  \frac{1}{4x^2} - \frac{2}{x} + \frac{17}{4} - x\]

 \[0 = 1 - 8x + 17x^2 - 4x^3.\]

 Then by the rational root theorem, this has roots $\frac{1}{4}$, $2 - \sqrt{3}$, and $2 + \sqrt{3}$. The first and last roots are extraneous because they imply $a = 0$ and $x > 1$, respectively, thus $x = \boxed{\textbf{(C)}\ 2-\sqrt{3}}$.

 



\textbf{Solution 4}

Let $\angle FKH$ = $k$ and $CF$ = $a$.  It is shown that all four triangles in the picture are similar.  From the square side lengths:

 
 \[a + \sin(k) \cdot 1 + \cos(k) \cdot a = 1\]

 \[\sin(k)a + \cos(k) = 1\]Solving for $a$ we get:

 
 \[a = \frac{1-\sin(k)}{\cos(k) + 1} = \frac{1 - \cos(k)}{\sin(k)}\]

 \[(1-\sin(k)) \cdot sin(k) = (1 - \cos(k))\cdot(\cos(k) + 1)\]

 \[\sin(k)-\sin(k)^2 = \cos(k) + 1 - \cos(k)^2 - \cos(k)\]

 \[\sin(k)-\sin(k)^2 = \sin(k)^2\]

 \[1-\sin(k) = \sin(k)\]

 \[\sin(k) = \frac{1}{2}, \cos(k) = \frac{\sqrt 3}{2}\]

 \[a = \frac{1 - \frac{\sqrt 3}{2}}{\frac{1}{2}} = 2 - \sqrt 3\]




\textbf{Solution 5}

Note that $HJ$ is a diagonal of $JKHG$, so it must be equal in length to $FB$. Therefore, quadrilateral $FHJB$ has $FH\parallel BJ$, and $FB=HJ$, so it must be either an isosceles trapezoid or a parallelogram. But due to the slope of $FB$ and $HJ$, we see that it must be a parallelogram. Therefore, $FH=BJ$. But by the symmetry in rectangle $FEAD$, we see that $FH=JA$. Therefore, $BJ=FH=JA$. We also know that $BJ+JA=1$, hence $BJ=JA=\frac12$.

 As $JG=1$ and $JA=\frac12$, and as $\triangle GJA$ is right, we know that $\triangle GJA$ must be a 30-60-90 triangle. Therefore, $GA=\sqrt{3}/2$ and $DG=1-\sqrt{3}/2$. But by similarity, $\triangle DHG$ is also a 30-60-90 triangle, hence $DH=\sqrt{3}-3/2$. But $\triangle DHG\cong\triangle EJK$, hence $EJ=\sqrt{3}-3/2$. As $BJ=1/2$, this implies that $BE=BJ-EJ=1/2-\sqrt{3}+3/2=2-\sqrt{3}$. Thus the answer is $\boxed{\textbf{(C)}}$.

 



\textbf{Solution 6}

Let $BE = x$, $KE = y$
$DG = y$, $AG = 1-y$, $KJ = x$, $JG = 1$
$\frac{AJ}{JG} = \frac{KE}{JK}$, $AJ = \frac{y}{x}$, $EJ = 1 - x - \frac{y}{x}$, 

 $\frac{AG}{JG} = \frac{EJ}{JK}$, $1-y = \frac{1 - x - \frac{y}{x}}{x}$, $x - xy = 1 - x - \frac{y}{x}$, $2x^2 - x^2y + y - x = 0$ $(1)$
$AJ^2 + AG^2 = JG^2$, $\left( \frac{y}{x} \right) ^2 + (1-y)^2 = 1$, $\frac{y^2}{x^2} + 1 - 2y + y^2 = 1$, $x^2 y -2x^2 + y = 0$ $(2)$
$(1) + (2)$, $2y - x = 0$, $y = \frac{x}{2}$

 Substitute $y$ into $(2)$ we get $x^2 \cdot \frac{x}{2} - 2 x^2 + \frac{x}{2} = 0$, $x^2 - 4x + 1 = 0$
$x = \frac{4 \pm \sqrt{16-4}}{2} = 2 \pm \sqrt{3}$, as $x < 1$, $x = \boxed{\textbf{(C)}2 - \sqrt{3}}$



	\item (\textit{AMC 12 12B 2014}) In a small pond there are eleven lily pads in a row labeled 0 through 10.  A frog is sitting on pad 1.  When the frog is on pad $N$, $0<N<10$, it will jump to pad $N-1$ with probability $\frac{N}{10}$ and to pad $N+1$ with probability $1-\frac{N}{10}$.  Each jump is independent of the previous jumps.  If the frog reaches pad 0 it will be eaten by a patiently waiting snake.  If the frog reaches pad 10 it will exit the pond, never to return.  What is the probability that the frog will escape without being eaten by the snake?


 $\textbf{(A) }\frac{32}{79}\qquad \textbf{(B) }\frac{161}{384}\qquad \textbf{(C) }\frac{63}{146}\qquad \textbf{(D) }\frac{7}{16}\qquad \textbf{(E) }\frac{1}{2}\qquad$

 \textbf{Answer: }C



\textbf{Solution 1}

A long, but straightforward bash:

 Define $P(N)$ to be the probability that the frog survives starting from pad N.

 Then note that by symmetry, $P(5) = 1/2$, since the probabilities of the frog moving subsequently in either direction from pad 5 are equal.

 We therefore seek to rewrite $P(1)$ in terms of $P(5)$, using the fact that

 $P(N) = \frac {N} {10}P(N - 1) + \frac {10 - N} {10}P(N + 1)$

 as said in the problem.

 Hence $P(1) = \frac {1} {10}P(0) + \frac {9} {10}P(2) = \frac {9} {10}P(2)$
$\Rightarrow P(2) = \frac {10} {9}P(1)$

 Returning to our original equation:

 $P(1) = \frac {9} {10}P(2) = \frac {9} {10}\left(\frac{2} {10}P(1) + \frac{8} {10}P(3)\right)$
$= \frac {9} {50}P(1) + \frac {18} {25}P(3) \Rightarrow P(1) - \frac {9} {50}P(1)$ $= \frac {18} {25}P(3)$
$\Rightarrow P(3) = \frac {41} {36}P(1)$

 Returning to our original equation:

 $P(1) = \frac {9} {50}P(1) + \frac {18} {25}\left(\frac {3} {10}P(2) + \frac {7} {10}P(4)\right)$
$= \frac {9} {50}P(1) + \frac {27} {125}P(2) + \frac {63} {125}P(4)$
$= \frac {9} {50}P(1) + \frac {27} {125}\left(\frac {10} {9}P(1)\right) + \frac {63} {125}\left(\frac {4} {10}P(3) + \frac {6} {10}P(5)\right)$

 Cleaning up the coefficients, we have: 

 $= \frac {21} {50}P(1) + \frac {126} {625}P(3) + \frac {189} {625}P(5)$
$= \frac {21} {50}P(1) + \frac {126} {625}\left(\frac {41} {36}P(1)\right) + \frac {189} {625}\left(\frac {1} {2}\right)$

 Hence, $P(1) = \frac {525} {1250}P(1) + \frac {287} {1250}P(1) + \frac {189} {1250}$
$\Rightarrow P(1) - \frac {812} {1250}P(1) = \frac {189} {1250} \Rightarrow P(1) = \frac {189} {438}$
$= \boxed{\frac {63} {146}\, (C)}$

 Or set $P(1)=a,P(2)=b,P(3)=c,P(4)=d,P(5)=e=1/2$:
 \[a=0.1\emptyset+0.9b,b=0.2a+0.8c,c=0.3b+0.7d,d=0.4c+0.6e\]
 \[10a=\emptyset+9b,10b=2a+8c,10c=3b+7d,10d=4c+6e\]
 \[\implies b=\frac{10a-\emptyset}{9},c=\frac{5b-a}{4},d=\frac{10c-3b}{7},e=\frac{5d-2c}{3}=1/2\]$b=\frac{10a}{9}$
$c=\frac{5\left(\frac{10a}{9}\right)-a}{4}=\frac{\frac{50a}{9}-a}{9}=\frac{41a}{36}$
$d=\frac{10\left(\frac{41a}{36}\right)-3\left(\frac{30a}{9}\right)}{7}=\frac{\frac{205a}{18}-\frac{10a}{3}}{7}=\frac{145a}{126}$
$e=\frac{5\left(\frac{145a}{126}\right)-2\left(\frac{41a}{36}\right)}{3}=\frac{\frac{725a}{126}-\frac{41a}{18}}{3}=\frac{73a}{63}$

 Since $e=\frac{1}{2}$, $\frac{73a}{63}=\frac{1}{2}\implies a=\boxed{\textbf{(C) }\frac{63}{146}}$.

 



\textbf{Solution 2}

Notice that the probabilities are symmetrical around the fifth lily pad. If the frog is on the fifth lily pad, there is a $\frac{1}{2}$ chance that it escapes and a $\frac{1}{2}$ chance that it gets eaten. Now, let $P_k$ represent the probability that the frog escapes if it is currently on pad $k$. We get the following system of $5$ equations:
 \[P_1=\frac{9}{10}\cdot P_2\]
 \[P_2=\frac{2}{10}\cdot P_1 + \frac{8}{10}\cdot P_3\]
 \[P_3=\frac{3}{10}\cdot P_2 + \frac{7}{10}\cdot P_4\]
 \[P_4=\frac{4}{10}\cdot P_3 + \frac{6}{10}\cdot P_5\]
 \[P_5=\frac{5}{10}\]We want to find $P_1$, since the frog starts at pad $1$. Solving the above system yields $P_1=\frac{63}{146}$, so the answer is $\boxed{(C)}$.

 System of equations (alligator112): Start by removing all fractions. From here on out, we will use substitution while avoiding fractions.
 \[10P_1 = 9P_2\]
 \[10P_2 = 2P_1 + 8P_3\]
 \[10P_3 = 3P_2 + 7P_4\]
 \[10P_4 = 4P_3 + 6P_5 = 4P_3 + 3\]
 \[100 P_3 = 30P_2 + 70P_4 = 30P_2 + 7(4P_3 + 3) \rightarrow 72P_3 = 30P_2 + 21\]
 \[90P_2 = 18P_1 + 72P_3 = 18P_1 + 30P_2 + 21 \rightarrow 60P_2 = 18P_1 + 21\]
 \[200P_1 = 180P_2 = 3(18P_1 + 21) = 54P_1 + 63 \rightarrow P_1 = \frac{63}{146}\]




\textbf{Solution 3}

Assign each lily pad a value, with pad $0$ having value $0$ and pad $i > 0$ having value $\sum_{k=0}^{i-1}\frac{1}{\binom{9}{k}}$. If we treat the process as a random walk $X_0, X_1, \cdots$ over these values, note that $\mathbb{E}[X_{n+1}] = n$ for all $n \ge 0$. This is purely by construction, as the expected gain at pad $k$ is $\binom{9}{k}\cdot \frac{10-k}{10} - \binom{9}{k-1}\cdot \frac{k}{10} = 0$. Therefore the process is a martingale. Since the stopping time has finite expectation and the increments are bounded, the \href{https://en.wikipedia.org/wiki/Optional\_stopping\_theorem}{Optional Stopping Thoerem} applies and the expected final outcome is $0$. This gives the probability of reaching safety as $\left(\sum_{k=0}^{9}\frac{1}{\binom{9}{k}}\right)^{-1} = \frac{63}{146}$, which is choice $\boxed{(C)}$.

 


 



\textbf{Solution 4 (Explaining the Martingale Shortcut Mentioned Previously)}

Instead of grinding through a massive system of 9 equations, we can solve this problem by assigning a "score value" to each lily pad. 

 Normally, if a game is perfectly fair (like a 50/50 coin flip), your average score after a turn is the same as your score before the turn. But this frog pond is rigged because the probabilities change at every pad. To fix this, we can mathematically design a custom scoring system that forces the game to be perfectly fair. 

 Once the scoring system is perfectly balanced, the law of expected value tells us that the probability of winning is just a simple ratio:
 \[\text{Escape Probability} = \frac{\text{Score of your Starting Pad}}{\text{Score of the Winning Pad (Pad 10)}}\]

            \begin{itemize}
        \item (Note: We set the score of Pad 0 to be 0 points, so getting eaten by the snake gives 0 credit).*
            \end{itemize}
            \textbf{Step 1: Find the Step Rule} To keep our scores perfectly balanced at any pad $N$, the reward for a rare forward jump must perfectly offset the risk of an easy backward jump. This gives us a simple rule for how the point-gaps between pads must scale:

 
 \[\text{Step Multiplier} = \frac{\text{Probability of jumping Left}}{\text{Probability of jumping Right}} = \frac{\frac{N}{10}}{\frac{10-N}{10}} = \frac{N}{10-N}\]

 This fraction $\frac{N}{10-N}$ tells us exactly how much larger the next point-step must be to keep our scoring scale balanced.

 \textbf{Step 2: Build the Chain of Gaps} We start by declaring that moving from Pad 0 to Pad 1 awards exactly $1$ point. To find the point value of each subsequent step, we take our current step value and multiply it by our Step Multiplier rule:

 
            \begin{itemize}
        \item \textbf{Step 1 (Pad 0 to 1):} Worth $1$ point. Pad 1 Score = $1$.
        \item \textbf{Step 2 (Pad 1 to 2):} Multiply by the rule at $N=1 \implies 1 \times \frac{1}{9} = \frac{1}{9}$ points.
        \item \textbf{Step 3 (Pad 2 to 3):} Multiply by the rule at $N=2 \implies \frac{1}{9} \times \frac{2}{8} = \frac{1}{36}$ points.
        \item \textbf{Step 4 (Pad 3 to 4):} Multiply by the rule at $N=3 \implies \frac{1}{36} \times \frac{3}{7} = \frac{1}{84}$ points.
        \item \textbf{Step 5 (Pad 4 to 5):} Multiply by the rule at $N=4 \implies \frac{1}{84} \times \frac{4}{6} = \frac{1}{126}$ points.
            \end{itemize}
            Because the pond's layout is perfectly symmetric, the point values for the remaining five steps mirror these numbers exactly in reverse order. This gives us all 10 step values:
 \[\left[ 1, \, \frac{1}{9}, \, \frac{1}{36}, \, \frac{1}{84}, \, \frac{1}{126}, \, \frac{1}{126}, \, \frac{1}{84}, \, \frac{1}{36}, \, \frac{1}{9}, \, 1 \right]\]

 \textbf{Step 3: Find the Total Score and Solve} Now we can easily calculate our two target scores:

 
            \begin{enumerate}
        \item \textbf{Score of our Starting Pad (Pad 1):} This is just the value of the very first step: $1$.
        \item \textbf{Score of the Winning Pad (Pad 10):} This is the total sum of all 10 steps combined. Grouping the identical matching pairs over a common denominator of 126 makes this simple to calculate:
            \end{enumerate}
            
 \[\text{Total Score of Pad 10} = 2 \left(1 + \frac{1}{9} + \frac{1}{36} + \frac{1}{84} + \frac{1}{126}\right) = 2 \left(\frac{252 + 28 + 7 + 3 + 2}{126}\right) = \frac{146}{63}\]

 Finally, we divide our starting score by the winning pad's score to get the final probability:

 
 \[P = \frac{\text{Starting Score}}{\text{Winning Score}} = \frac{1}{\frac{146}{63}} = \boxed{\textbf{(C) } \frac{63}{146}}\]

 \textbf{Remark (Markov Chain)} 
            \begin{enumerate}[label=(\roman*)]
                \item <i>Main article: \href{https://artofproblemsolving.com/wiki/index.php?title=Markov\_Chains}{Markov Chains}</i>
            \end{enumerate}
            We can represent this problem with the following State Transition Diagram of \href{https://en.wikipedia.org/wiki/Absorbing\_Markov\_chain}{Absorbing Markov Chain}. State $0$ and state $10$ are the absorbing states. This problem asks for the \href{https://en.wikipedia.org/wiki/Absorbing\_Markov\_chain\#Absorbing\_probabilities}{Absorbing Probability} of state $10$.

 
                    \begin{figure}[H]
                        \centering
                        \includegraphics[width=0.5\textwidth]{Solution/AMC_12_12B_2014_22_Solution_Figure_1.png}
                    \end{figure}
                    Let $p_{ij} = P(X_{n+1} = j | X_n = i)$, the probability that state $i$ transits to state $j$ on the next step. 

 Let $a_i$ be the absorbing probability of being absorbed in the absorbing state when starting from transient state $i$.

 $a_i = \sum_{j} (p_{ij} \cdot a_{j})$
$a_i$ is the sum of the products of $p_{ij}$ and $a_j$ of all the next state $j$.

 Note:

 This sum results in the same system of equations as solution 2, and the solution follows as well. 

 


	\item (\textit{AMC 12 12B 2014}) The number 2017 is prime.  Let $S = \sum \limits_{k=0}^{62} \dbinom{2014}{k}$.  What is the remainder when $S$ is divided by 2017?


 $\textbf{(A) }32\qquad \textbf{(B) }684\qquad \textbf{(C) }1024\qquad \textbf{(D) }1576\qquad \textbf{(E) }2016\qquad$

 \textbf{Answer: }C



\textbf{Solution 1}

Note that $2014\equiv -3 \mod2017$.  We have for $k\ge1$
 \[\dbinom{2014}{k}\equiv \frac{(-3)(-4)(-5)....(-2-k)}{k!}\mod 2017\] 
 \[\equiv (-1)^k\dbinom{k+2}{k} \mod 2017\]Therefore 
 \[\sum \limits_{k=0}^{62} \dbinom{2014}{k}\equiv \sum \limits_{k=0}^{62}(-1)^k\dbinom{k+2}{2} \mod 2017\]This is simply an alternating series of triangular numbers that goes like this: $1-3+6-10+15-21....$After finding the first few sums of the series, it becomes apparent that 
 \[\sum \limits_{k=1}^{n}(-1)^k\dbinom{k+2}{2}\equiv -\left(\frac{n+1}{2} \right) \left(\frac{n+1}{2}+1 \right) \mod 2017 \textnormal{  if n is odd}\]and 
 \[\sum \limits_{k=1}^{n}(-1)^k\dbinom{k+2}{2}\equiv \left(\frac{n}{2}+1 \right)^2 \mod 2017 \textnormal{  if n is even}\]Obviously, $62$ falls in the second category, so our desired value is 
 \[\left(\frac{62}{2}+1 \right)^2 = 32^2 = \boxed{\textbf{(C)}\ 1024}\]

 \textbf{Sidenote} Another way to finish, using the fact that $\dbinom{k+2}{2} = 1 + 2 + \dots + (k+1)$:
 \begin{align*} \sum \limits_{k=0}^{62}(-1)^k\dbinom{k+2}{2} &\equiv \sum \limits_{k=1}^{63}(-1)^{k-1} (1 + 2 + \dots + k) \\ &\equiv 1 - (1+2) + (1+2+3) - (1+2+3+4) + \dots + (1 + \dots + 63) \\ &\equiv 1 + 3 + 5 + \dots + 63 \\ &\equiv \boxed{1024} \mod 2017 \end{align*}





\textbf{Solution 2 (By Junfei)}

We want to find the value of an extremely complicated sum$\pmod{2017}$. The key idea for such problems is to use various methods to convert numbers in the form $\dbinom{2014}{i}$ to numbers in the form $\dbinom{2017}{\text{some integer}}$. Considering the fact that $2014$ and $2017$ are very close numbers, we might be able to apply Pascal's Theorem, which states that
 \[\dbinom{2014}{i}+\dbinom{2014}{i+1}=\dbinom{2015}{i+1}\].

 
By applying Pascal's Theorem $3$ times, we can prove the identity
 \[\dbinom{2017}{i+3}=\dbinom{2014}{i}+3\dbinom{2014}{i+1}+3\dbinom{2014}{i+2}+\dbinom{2014}{i+3}\]

 
The path is clear: we need to design an ingenious method to boil down the large numbers in the sum, such as $\dbinom{2014}{62}$, which is impossible to calculate. We need to express these numbers using sums and differences of binomial coefficients with $2017$ as the number of elements.

 
After some trial, we may conceive that doing vertical calculation in a way similar to those of long divisions (see picture), where we keep adding and subtracting binomial coefficients with $2017$ as the number of elements to cancel out every new term with $2014$ as the number of elements, might provide us with some clues.

 


 
                    \begin{figure}[H]
                        \centering
                        \includegraphics[width=0.5\textwidth]{Solution/AMC_12_12B_2014_23_Solution_Figure_1.png}
                    \end{figure}
                    We can see that, to express a single term such as $\dbinom{2014}{1}$, the multiples of $\dbinom{2017}{4},\dbinom{2017}{5},\dbinom{2017}{6},\dbinom{2017}{7},\dbinom{2017}{8}......$ we require are

 
 \[1,-3,6,-10,15......\]

 If we have adequate number sense, we can wisely notice that these terms are the first terms of the sequence $A_n=\frac{n^2+n}{2}$·$(-1)^{n+1}$, where the terms with even $n$ are turned negative.

 
Expressing a term in the form $\dbinom{2014}{i}$ using this sequence of multiples of $\dbinom{2017}{\text{some integer}}$, we can see that all the terms before $\dbinom{2017}{2017}$ are multiples of $2017$ and can therefore be neglected. The only term we need to consider is the final term, $A_{2015-i}\dbinom{2017}{2017}=A_{2015-i}$.

 
By writing down the expression of $A_n$ and simplifying, we can routinely verify the $A_n+3A_{n+1}+3A_{n+2}+A_{n+3}=0$ identity that we have discovered.

 
For the expansion of the right-hand side of the expression
 \[\dbinom{2014}{i}=\dbinom{2017}{i+3}+A_2{\dbinom{2017}{i+4}}+A_3{\dbinom{2017}{i+5}}+A_4{\dbinom{2017}{i+6}}+......+A_{2015-i}{\dbinom{2017}{2017}}\],all binomial coefficients with $2014$ as the number of elements are completely cancelled out except for the term $\dbinom{2014}{i}$ itself, which appears exactly once——it is contributed by the term $\dbinom{2017}{i+3}$.

 
In conclusion,

 
 \[\dbinom{2014}{i}=A_{2015-i}\pmod{2017}\]

 Now we only have to find the remainder of $1+\sum_{k=1953}^{2014}A_k$ upon division by $2017$. The remainder of the terms of $A_n$ upon division by $2017$ can be calculated quickly using modular arithmetic.

 
The result of the expression is $1-1954-1956-1958-......-2014$
$=1+63+61+59+......+3$
$=1024\pmod{2017}$.

 
$1024$ is the answer to the problem.

 


	\item (\textit{AMC 12 12B 2014}) Let $ABCDE$ be a pentagon inscribed in a circle such that $AB = CD = 3$, $BC = DE = 10$, and $AE= 14$.  The sum of the lengths of all diagonals of $ABCDE$ is equal to $\frac{m}{n}$, where $m$ and $n$ are relatively prime positive integers.  What is $m+n$ ?


 $\textbf{(A) }129\qquad \textbf{(B) }247\qquad \textbf{(C) }353\qquad \textbf{(D) }391\qquad \textbf{(E) }421\qquad$

 \textbf{Answer: }D



\textbf{Solution 1}

Let $BE=a$, $AD=b$, and $AC=CE=BD=c$. Let $F$ be on $AE$ such that $CF \perp AE$. 
                    \begin{figure}[H]
                        \centering
                        \includegraphics[width=0.5\textwidth]{Solution/AMC_12_12B_2014_24_Solution_Figure_1.png}
                    \end{figure}
                    In $\triangle CFE$ we have $\cos\theta = -\cos(\pi-\theta)=-7/c$. We use the \href{https://artofproblemsolving.com/wiki/index.php?title=Law\_of\_Cosines}{Law of Cosines} on $\triangle ABC$ to get $60\cos\theta = 109-c^2$.  Eliminating $\cos\theta$ we get $c^3-109c-420=0$ which factorizes as 
 \[(c+7)(c+5)(c-12)=0.\]Discarding the negative roots we have $c=12$. Thus $BD=AC=CE=12$. For $BE=a$, we use Ptolemy's theorem on cyclic quadrilateral $ABCE$ to get $a=44/3$. For $AD=b$, we use \href{https://artofproblemsolving.com/wiki/index.php?title=Ptolemy\%27s\_theorem}{Ptolemy's theorem} on cyclic quadrilateral $ACDE$ to get $b=27/2$.

 The sum of the lengths of the diagonals is $12+12+12+\frac{44}{3}+\frac{27}{2} = \frac{385}{6}$ so the answer is $385 + 6 = \fbox{\textbf{(D) }391}$




\textbf{Solution 2}

Let $a$ denote the length of a diagonal opposite adjacent sides of length $14$ and $3$, $b$ for sides $14$ and $10$, and $c$ for sides $3$ and $10$. Using Ptolemy's Theorem on the five possible quadrilaterals in the configuration, we obtain:

 
 \begin{align} c^2 &= 3a+100 \\ c^2 &= 10b+9 \\ ab &= 30+14c \\ ac &= 3c+140\\ bc &= 10c+42 \end{align}

 Using equations $(1)$ and $(2)$, we obtain:

 
 \[a = \frac{c^2-100}{3}\]

 and

 
 \[b = \frac{c^2-9}{10}\]

 Plugging into equation $(4)$, we find that:

 
 \begin{align*} \frac{c^2-100}{3}c &= 3c + 140\\ \frac{c^3-100c}{3} &= 3c + 140\\ c^3-100c &= 9c + 420\\ c^3-109c-420 &=0\\ (c-12)(c+7)(c+5)&=0 \end{align*}


 Or similarly into equation $(5)$ to check:

 
 \begin{align*} \frac{c^2-9}{10}c &= 10c+42\\ \frac{c^3-9c}{10} &= 10c + 42\\ c^3-9c &= 100c + 420\\ c^3-109c-420 &=0\\ (c-12)(c+7)(c+5)&=0 \end{align*}

$c$, being a length, must be positive, implying that $c=12$. In fact, this is reasonable, since $10+3\approx 12$ in the pentagon with apparently obtuse angles. Plugging this back into equations $(1)$ and $(2)$ we find that $a = \frac{44}{3}$ and $b= \frac{135}{10}=\frac{27}{2}$.

 We desire $3c+a+b = 3\cdot 12 + \frac{44}{3} + \frac{27}{2} = \frac{216+88+81}{6}=\frac{385}{6}$, so it follows that the answer is $385 + 6 = \fbox{\textbf{(D) }391}$




\textbf{Solution 3 (Ptolemy's but Quicker)}

Let us set $x$ to be $AC=BD=CE$ and $y$ to be $BE$ and $z$ to be $AD$. It follow from applying \href{https://artofproblemsolving.com/wiki/index.php?title=Ptolemy\%27s\_Theorem}{Ptolemy's Theorem} on $ABCD$ to get $x^2=9+10z$. Applying Ptolemy's on $ACDE$ gives $xz=42+10x$; and applying Ptolemy's on $BCDE$ gives $x^2=100+3y$. So, we have the have the following system of equations:

 
 \begin{align} x^2 &= 9+10z \\ x^2 &= 100+3y \\ xz &= 42+10x \end{align}

 From $(3)$, we have $42=(z-10)x$. Isolating the x gives $x=\dfrac{42}{z-10}$. Plugging in $x=\dfrac{42}{z-10}$ to $(1)$ gives

 
 \begin{align*} \left(\frac{42}{z-10}\right)^2 &= 10z+9\\ 10z^3 - 191z^2 + 820z + 900 &= 1764\\ 10z^3 - 191z^2 + 820z - 864 &= 0\\ (5z-8)(2z-27)(z-4) &=0 \end{align*}


 It is impossible for $z<10$ for $x<0$; that means $z=\frac{27}{2}$. That means $x = 12$ and $y = \frac{44}{3}$.

 Thus, the sum of all diagonals is $3x+y+z = 3\cdot 12 + \frac{44}{3} + \frac{27}{2} = 385/6$, which implies our answer is $m+n = 385+6 = \fbox{391 \textbf{(D)}}$.

 



\textbf{Solution 4}


                    \begin{figure}[H]
                        \centering
                        \includegraphics[width=0.5\textwidth]{Solution/AMC_12_12B_2014_24_Solution_Figure_2.png}
                    \end{figure}
                    Let $BE = a$, $AC = CE = BD = b$

 By \href{https://artofproblemsolving.com/wiki/index.php?title=Ptolemy\%27s\_theorem}{Ptolemy's theorem} for quadrilateral $ABCE$, $AB \cdot CE + BC \cdot AE = BE \cdot AC$, $3b + 140 = ab$, $a = 3 + \frac{140}{b}$

 By \href{https://artofproblemsolving.com/wiki/index.php?title=Ptolemy\%27s\_theorem}{Ptolemy's theorem} for quadrilateral $BCDE$, $CD \cdot BE + BC \cdot DE = BD \cdot CE$, $3a + 100 = b^2$
$3(3 + \frac{140}{b}) + 100 = b^2$, $b^3 - 109 b -420 = 0$, $(b-12)(b+7)(b+5) = 0$, $b = 12$
$a = 3 + \frac{140}{12} = \frac{44}{3}$

 By \href{https://artofproblemsolving.com/wiki/index.php?title=Ptolemy\%27s\_theorem}{Ptolemy's theorem} for quadrilateral $ABDE$, $AE \cdot BD + AB \cdot DE = AD \cdot BE$, $AD \cdot a = 14b + 30$, $AD = \frac{27}{2}$
$\frac{m}{n} = 12 + 12 + 12 + \frac{44}{3} + \frac{27}{2} = \frac{385}{6}$, $385 + 6 = \boxed{\textbf{(D) }391}$



	\item (\textit{AMC 12 12B 2014}) Find the sum of all the positive solutions of 
 \[2\cos(2x) \left(\cos(2x) - \cos\left( \frac{2014\pi^2}{x} \right)\right) = \cos(4x) - 1\]
$\textbf{(A)}\ \pi \qquad\textbf{(B)}\ 810\pi  \qquad\textbf{(C)}\ 1008\pi \qquad\textbf{(D)}\ 1080 \pi \qquad\textbf{(E)}\ 1800\pi$

 \textbf{Answer: }D



\textbf{Solution 1}

Rewrite $\cos{4x} - 1$ as $2\cos^2{2x} - 2$.  Now let $a = \cos{2x}$, and let $b = \cos{\left( \frac{2014\pi^2}{x} \right) }$.  We have:
 \[2a(a - b) = 2a^2 - 2\]

 Therefore, 
 \[ab = 1\].

 Notice that either $a = 1$ and $b = 1$ or $a = -1$ and $b = -1$.  For the first case, $a = 1$ only when $x = k\pi$ and $k$ is an integer.  $b = 1$ when $\frac{2014\pi^2}{k\pi}$ is an even multiple of $\pi$, and since $2014 = 2*19*53$, $b =1$ only when $k$ is an odd divisor of $2014$.  This gives us these possible values for $x$:
 \[x= \pi, 19\pi, 53\pi, 1007\pi\]For the case where $a = -1$, $\cos{2x} = -1$, so $x = \frac{m\pi}{2}$, where m is odd.  $\frac{2014\pi^2}{\frac{m\pi}{2}}$ must also be an odd multiple of $\pi$ in order for $b$ to equal $-1$, so $\frac{4028}{m}$ must be odd.  We can quickly see that dividing an even number by an odd number will never yield an odd number, so there are no possible values for $m$, and therefore no cases where $a = -1$ and $b = -1$. Therefore, the sum of all our possible values for $x$ is
 \[\pi + 19\pi + 53\pi + 1007\pi = \boxed{\textbf{(D)}\ 1080 \pi}\]




\textbf{Solution 2}

Very similar to the solution above, re-write the expression using $\cos 4x = 2 \cos^2 2x - 1$:

 
 \[2\cos2x \left(\cos2x - \cos{\left( \frac{2014\pi^2}{x} \right) } \right) = 2 \cos^2 2x - 2\]

 Now, expand the LHS and cancel terms:

 
 \[2 \cos^2 2x - 2 \cos 2x \cos{\left( \frac{2014\pi^2}{x} \right) } = 2 \cos^2 2x - 2\]

 \[\cos 2x \cos{\left( \frac{2014\pi^2}{x} \right) } = 1\]

 Now we use product-to-sum identities to get:

 
 \[\frac{1}{2} \Biggl( \cos\Bigl( 2x + \frac{2014 \pi^2}{x} \Bigr) + \cos\Bigl( 2x + \frac{2014 \pi^2}{x} \Bigr) \Biggr) = 1\]

 \[\cos{\left(2x + \left( \frac{2014\pi^2}{x} \right) \right)} = 1\]

 Notice that for any $\theta$, $\max{\cos \theta} = 1$. This is achieved when $\theta = 0$, or equivalently 

 
 \[2x - \frac{2014\pi^2}{x} \equiv 0 \mod 2\pi\]

 We can cleverly assume $x=c\pi$ for some real $c$. Then, we must have

 
 \[2c - \frac{2014}{c} \equiv 0 \mod 2\]

 In order for this to be satisfied, $\frac{2014}{c}$ must be an even integer. Factoring $2014 = 2 \cdot 19 \cdot 53$, we see that $c$ must be a positive odd integer that divides $2014$. Our only positive valid $c$ are $c = 1, 19, 53, 1007$. Our answer is just $\pi(1+19+53+1007) = 1080\pi \implies \textbf{(D)}$. 

 
-FIREDRAGONMATH16

 



\textbf{Solution 3}

Rewriting $\cos{4x} - 1$ as $2\cos^2{2x} - 2$ and transposing $2\cos{2x}$ from the LHS to the RHS, we get,

 
 \[\cos{2x} - \cos{\left(\frac{2014\pi^2}{x}\right)} = 1 - \frac{1}{\cos{2x}}\]
 \[\implies \underbrace{\cos{2x} + \frac{1}{\cos{2x}}}_{\text{LHS}} = \underbrace{1 + \cos{\left(\frac{2014\pi^2}{x}\right)}}_{\text{RHS}}\]

 By the \href{https://artofproblemsolving.com/wiki/index.php/AM-GM\_Inequality}{AM-GM Inequality},

 
 \[\cos{2x} + \frac{1}{\cos{2x}} \in (-\infty, -2] \cup [2, \infty)\]

 Also, because of the range of $\cos$,

 
 \[1 + \cos{\left(\frac{2014\pi^2}{x}\right)} \in [0, 2]\]

 Hence, $\text{LHS} = \text{RHS} = 2$, and we get ($m, n \in \mathbb{Z}$),

 
            \begin{itemize}
        \item $\cos{2x} = 1 \implies x = m\pi \hspace{75pt} \cdots\cdots (1)$
        \item $\cos{\left(\frac{2014\pi^2}{x}\right)} = 1 \implies x = \frac{1007\pi}{n} \hspace{22pt} \cdots\cdots (2)$
            \end{itemize}
            
From $(1)$ and $(2)$,
 \[m = \frac{1007}{n}\]
 \[\implies m \in \{1, 19, 53, 1007\}\]
 \[\implies x \in \{1\pi, 19\pi, 53\pi, 1007\pi\}\]

 Therefore, sum of values of $x$ is 
 \[\pi + 19\pi + 53\pi + 1007\pi = \boxed{\textbf{(D)}\ 1080\pi}\].

 


\end{enumerate}